% GENERATED FILE. Edit canonical lessons, then run npm run generate:books.
% canonical-source-hash: fnv1a64:69079fc7ca91e86b
% canonical-lessons: BN-C459-binoyi, BN-C459-nomro, BN-C459-uddhoto, BN-C459-beyadob, BN-C459-toshamod-kora

\chapter{Modest and Rude}
\label{ch:bn-modest-and-rude}
\hlchaptermodality{459}

\begin{chapteropening}
\textbf{By the end of this chapter:} \emph{I can say whether someone is modest, gentle or rude.}

Complete the last lesson of chapter 459: I can say whether someone is modest, gentle or rude.
\end{chapteropening}

\section[\texorpdfstring{\bn{বিনয়ী} (binôyī)}{বিনয়ী (binôyī)}]{\bn{বিনয়ী} (binôyī) — modest; humble}

\label{lesson:BN-C459-binoyi}



\begin{quote}
Before the new one: say the Bengali for cunning, then the Bengali for to cheat.
\end{quote}

\subsection*{You'll want to know: \bn{বিনয়ী}}
\textbf{\bn{বিনয়ী}} — \emph{binôyī} — \textquotedblleft{}modest; humble\textquotedblright{}.

\begin{quote}
\bn{গায়িকা} \bn{খুব} \bn{বিখ্যাত}, \bn{কিন্তু} \bn{তিনি} \bn{বিনয়ী।}
\end{quote}

\begin{quote}
\emph{gāyikā khub bikhyātô, kintu tini binôyī.}
\end{quote}

\begin{quote}
The singer is very famous, but she is modest.
\end{quote}

\begin{cousinweb}
\textbf{\bn{বিনয়ী}} (\emph{binôyī}) is \emph{binôy}, Sanskrit \emph{vinaya}, \textquotedblleft{}training, discipline, modesty\textquotedblright{} (Monier-Williams), with the ending \emph{-ī}, \textquotedblleft{}having\textquotedblright{}.
\end{cousinweb}

\begin{grammarlens}[title={What you've built}]
One more word for describing what you lived and felt.
\end{grammarlens}

\subsection*{Your turn}
\begin{itemize}
\raggedright
  \item \emph{Say it:} \emph{binôyī}
  \item \emph{Say it:} \emph{binôyī}, once more
  \item \emph{Say it:} \emph{gāyikā khub bikhyātô, kintu tini binôyī.} — The singer is very famous, but she is modest.
  \item \emph{Recall:} say the Bengali for cunning, then the Bengali for to cheat, then say \emph{binôyī} again
\end{itemize}

\begin{tcolorbox}[breakable,colback=teal!4,colframe=teal!35!black,title={Before you move on}]
What does \bn{বিনয়ী} mean? (\textquotedblleft{}Modest; humble\textquotedblright{}.) Say it once more.
\end{tcolorbox}
\section[\texorpdfstring{\bn{নম্র} (nômrô)}{নম্র (nômrô)}]{\bn{নম্র} (nômrô) — gentle; soft-spoken}

\label{lesson:BN-C459-nomro}



\begin{quote}
Before the new one: say the Bengali for to cheat, then the Bengali for modest.
\end{quote}

\subsection*{You'll want to know: \bn{নম্র}}
\textbf{\bn{নম্র}} — \emph{nômrô} — \textquotedblleft{}gentle; soft-spoken\textquotedblright{}.

\begin{quote}
\bn{নার্সটি} \bn{রোগীদের} \bn{সঙ্গে} \bn{নম্র।}
\end{quote}

\begin{quote}
\emph{nārsṭi rogīder sôngge nômrô.}
\end{quote}

\begin{quote}
The nurse is gentle with the patients.
\end{quote}

\begin{cousinweb}
\textbf{\bn{নম্র}} (\emph{nômrô}) is a tatsama, Sanskrit \emph{namra}, \textquotedblleft{}bowing, bent\textquotedblright{} and so \textquotedblleft{}humble\textquotedblright{} (Monier-Williams).
\end{cousinweb}

\begin{grammarlens}[title={What you've built}]
One more word for describing what you lived and felt.
\end{grammarlens}

\subsection*{Your turn}
\begin{itemize}
\raggedright
  \item \emph{Say it:} \emph{nômrô}
  \item \emph{Say it:} \emph{nômrô}, once more
  \item \emph{Say it:} \emph{nārsṭi rogīder sôngge nômrô.} — The nurse is gentle with the patients.
  \item \emph{Recall:} say the Bengali for to cheat, then the Bengali for modest, then say \emph{nômrô} again
\end{itemize}

\begin{tcolorbox}[breakable,colback=teal!4,colframe=teal!35!black,title={Before you move on}]
What does \bn{নম্র} mean? (\textquotedblleft{}Gentle; soft-spoken\textquotedblright{}.) Say it once more.
\end{tcolorbox}
\section[\texorpdfstring{\bn{উদ্ধত} (uddhôtô)}{উদ্ধত (uddhôtô)}]{\bn{উদ্ধত} (uddhôtô) — insolent; haughty}

\label{lesson:BN-C459-uddhoto}



\begin{quote}
Before the new one: say the Bengali for modest, then the Bengali for soft-spoken.
\end{quote}

\subsection*{You'll want to know: \bn{উদ্ধত}}
\textbf{\bn{উদ্ধত}} — \emph{uddhôtô} — \textquotedblleft{}insolent; haughty\textquotedblright{}.

\begin{quote}
\bn{উদ্ধত} \bn{ছেলেটা} \bn{দাদুর} \bn{সঙ্গে} \bn{তর্ক} \bn{করল।}
\end{quote}

\begin{quote}
\emph{uddhôtô chheleṭā dādur sôngge tôrkô kôrlô.}
\end{quote}

\begin{quote}
The insolent boy argued with my grandfather.
\end{quote}

\begin{cousinweb}
\textbf{\bn{উদ্ধত}} (\emph{uddhôtô}) is a tatsama, Sanskrit \emph{uddhata}, \textquotedblleft{}raised, lifted up\textquotedblright{}, and so \textquotedblleft{}haughty\textquotedblright{} (Monier-Williams): with the head held too high.
\end{cousinweb}

\begin{grammarlens}[title={What you've built}]
One more word for describing what you lived and felt.
\end{grammarlens}

\subsection*{Your turn}
\begin{itemize}
\raggedright
  \item \emph{Say it:} \emph{uddhôtô}
  \item \emph{Say it:} \emph{uddhôtô}, once more
  \item \emph{Say it:} \emph{uddhôtô chheleṭā dādur sôngge tôrkô kôrlô.} — The insolent boy argued with my grandfather.
  \item \emph{Recall:} say the Bengali for modest, then the Bengali for soft-spoken, then say \emph{uddhôtô} again
\end{itemize}

\begin{tcolorbox}[breakable,colback=teal!4,colframe=teal!35!black,title={Before you move on}]
What does \bn{উদ্ধত} mean? (\textquotedblleft{}Insolent; haughty\textquotedblright{}.) Say it once more.
\end{tcolorbox}
\section[\texorpdfstring{\bn{বেয়াদব} (beyādôb)}{বেয়াদব (beyādôb)}]{\bn{বেয়াদব} (beyādôb) — rude; ill-mannered (towards elders)}

\label{lesson:BN-C459-beyadob}



\begin{quote}
Before the new one: say the Bengali for soft-spoken, then the Bengali for insolent.
\end{quote}

\subsection*{You'll want to know: \bn{বেয়াদব}}
\textbf{\bn{বেয়াদব}} — \emph{beyādôb} — \textquotedblleft{}rude; ill-mannered (towards elders)\textquotedblright{}.

\begin{quote}
\bn{বাবা} \bn{বেয়াদব} \bn{ছেলেটাকে} \bn{বকলেন।}
\end{quote}

\begin{quote}
\emph{bābā beyādôb chheleṭāke bôklen.}
\end{quote}

\begin{quote}
Father scolded the ill-mannered boy.
\end{quote}

\begin{cousinweb}
\textbf{\bn{বেয়াদব}} (\emph{beyādôb}) is usually traced to Persian \emph{be-adab}, \textquotedblleft{}without manners\textquotedblright{}, where \emph{adab} is Arabic for good manners.
\end{cousinweb}

\begin{grammarlens}[title={What you've built}]
One more word for describing what you lived and felt.
\end{grammarlens}

\subsection*{Your turn}
\begin{itemize}
\raggedright
  \item \emph{Say it:} \emph{beyādôb}
  \item \emph{Say it:} \emph{beyādôb}, once more
  \item \emph{Say it:} \emph{bābā beyādôb chheleṭāke bôklen.} — Father scolded the ill-mannered boy.
  \item \emph{Recall:} say the Bengali for soft-spoken, then the Bengali for insolent, then say \emph{beyādôb} again
\end{itemize}

\begin{tcolorbox}[breakable,colback=teal!4,colframe=teal!35!black,title={Before you move on}]
What does \bn{বেয়াদব} mean? (\textquotedblleft{}Rude; ill-mannered (towards elders)\textquotedblright{}.) Say it once more.
\end{tcolorbox}
\section[\texorpdfstring{\bn{তোষামোদ} \bn{করা} (toshāmod kôrā)}{তোষামোদ করা (toshāmod kôrā)}]{\bn{তোষামোদ} \bn{করা} (toshāmod kôrā) — to flatter; to fawn on}

\label{lesson:BN-C459-toshamod-kora}



\begin{quote}
Before the new one: say the Bengali for insolent, then the Bengali for ill-mannered.
\end{quote}

\subsection*{You'll want to know: \bn{তোষামোদ} \bn{করা}}
\textbf{\bn{তোষামোদ} \bn{করা}} — \emph{toshāmod kôrā} — \textquotedblleft{}to flatter; to fawn on\textquotedblright{}.

\begin{quote}
\bn{চাকরির} \bn{জন্য} \bn{লোকটা} \bn{ম্যানেজারকে} \bn{তোষামোদ} \bn{করত।}
\end{quote}

\begin{quote}
\emph{chākrir jônyô lokṭā mænejārke toshāmod kôrtô.}
\end{quote}

\begin{quote}
The man used to flatter the manager to keep his job.
\end{quote}

Past (he/she …): \bn{তিনি} \bn{মন্ত্রীকে} \bn{তোষামোদ} \bn{করলেন}, \emph{tini môntrīke toshāmod kôrlen}.

\begin{cousinweb}
\textbf{\bn{তোষামোদ} \bn{করা}} (\emph{toshāmod kôrā}) is \emph{toshāmod}, \textquotedblleft{}flattery\textquotedblright{}, usually traced to Persian \emph{khushāmad}, with \emph{kôrā}, \textquotedblleft{}to do\textquotedblright{} (CDIAL 2814).
\end{cousinweb}

\begin{grammarlens}[title={What you've built}]
One more word for describing what you lived and felt.
\end{grammarlens}

\subsection*{Your turn}
\begin{itemize}
\raggedright
  \item \emph{Say it:} \emph{toshāmod kôrā}
  \item \emph{Say it:} \emph{toshāmod kôrā}, once more
  \item \emph{Say it:} \emph{chākrir jônyô lokṭā mænejārke toshāmod kôrtô.} — The man used to flatter the manager to keep his job.
  \item \emph{Recall:} say the Bengali for insolent, then the Bengali for ill-mannered, then say \emph{toshāmod kôrā} again
\end{itemize}

\begin{tcolorbox}[breakable,colback=teal!4,colframe=teal!35!black,title={Before you move on}]
What does \bn{তোষামোদ} \bn{করা} mean? (\textquotedblleft{}To flatter; to fawn on\textquotedblright{}.) Say it once more.
\end{tcolorbox}
